\documentclass[11pt,a4paper]{article}
\usepackage{rankit}
\usepackage[margin=2.4cm,headheight=16pt]{geometry}
\usepackage{fancyhdr}
\usepackage{microtype}
\usepackage{enumitem}

% Uzupełnij metrykę przed przygotowaniem nowego dokumentu.
\newcommand{\DocTitle}{[Tytuł dokumentu]}
\newcommand{\DocId}{[RPI-RANKIT-TYP-NNN]}
\newcommand{\DocVersion}{[1.0]}
\newcommand{\DocDate}{[RRRR-MM-DD]}
\newcommand{\DocAuthors}{[Imiona i nazwiska autorów]}
\newcommand{\DocStatus}{[Roboczy / do przeglądu / zatwierdzony]}

\pagestyle{fancy}
\fancyhf{}
\fancyhead[L]{\RankItWordmark\quad\small RPI 2026/2027}
\fancyhead[R]{\small \DocId}
\fancyfoot[C]{\small \DocId\quad|\quad v\DocVersion\quad|\quad\thepage}
\renewcommand{\headrulewidth}{0.4pt}
\setlength{\parindent}{0pt}
\setlength{\parskip}{0.5em}
\setlist{nosep,leftmargin=*}

\hypersetup{pdftitle={Szablon dokumentu projektu RankIt},pdfauthor={Zespół RankIt}}

\begin{document}
\begin{titlepage}
  \centering
  \vspace*{2cm}
  {\Huge\RankItWordmark\par}
  \vspace{0.8cm}
  {\Large Realizacja Projektu Informatycznego 2026/2027\par}
  \vspace{1.8cm}
  {\huge\bfseries\color{RankItNavy}\DocTitle\par}
  \vspace{0.6cm}
  {\color{RankItTeal}\rule{0.72\textwidth}{1.5pt}\par}
  \vspace{1.2cm}
  \begin{tabular}{rl}
    \textbf{Projekt:} & RankIt - aplikacja webowa do tworzenia rankingów \\
    \textbf{Identyfikator:} & \DocId \\
    \textbf{Wersja:} & \DocVersion \\
    \textbf{Data:} & \DocDate \\
    \textbf{Autorzy:} & \DocAuthors \\
    \textbf{Status:} & \DocStatus \\
    \textbf{Opiekun:} & [Imię i nazwisko] \\
  \end{tabular}
  \vfill
  {\small\url{https://github.com/elprofesoriqo/rpi}\par}
\end{titlepage}
\setcounter{page}{2}

\section{Historia zmian}
\begin{tabularx}{\textwidth}{p{0.14\textwidth}p{0.2\textwidth}p{0.25\textwidth}X}
  \toprule
  \textbf{Wersja} & \textbf{Data} & \textbf{Autor} & \textbf{Opis zmiany} \\
  \midrule
  1.0 & [RRRR-MM-DD] & [Autor] & Utworzenie dokumentu. \\
  \bottomrule
\end{tabularx}

\section{Cel i zakres}
Opisz, po co powstał dokument, do kogo jest skierowany i jakie informacje lub decyzje zawiera.

\section{Treść}
Podziel treść na krótkie sekcje z opisowymi nagłówkami. Oddziel fakty od planów i założeń. Przy każdej ważnej decyzji podaj właściciela albo źródło.

\section{Wnioski i dalsze kroki}
\begin{itemize}
  \item \emph{[Decyzja lub wniosek]}
  \item \emph{[Osoba odpowiedzialna i termin]}
\end{itemize}

\section{Źródła}
\begin{itemize}
  \item \emph{[Nazwa i adres źródła albo ścieżka do dokumentu]}
\end{itemize}

\end{document}
